\Needspace{7\baselineskip}
\begin{proof}[Proof of Proposition~\ref{prop:rewrite-uncapped-unit-bgp}]
\phantomsection
\label{proof:rewrite-uncapped-unit-bgp}
Use the definitions in Appendix~\ref{app:rewrite-uncapped-unit}.
\textit{The reference BGP.}
First, all candidate ratios are positive. The research denominator
$\rho-n+\eta g_Y^*$ is positive. Moreover,
\[
 i^*=\frac{\alpha(g_Y^*+\delta)}{\rho-n+g_Y^*+\delta}<\alpha,
 \qquad
 m^*<\frac{u^*\eta}{1-\eta}.
\]
Since $\Delta=(1-\alpha)(1-\eta-\omega_X)>0$,
$\omega_X<1-\eta$ and $\omega_X<1$. Consequently,
\[
 u^*+m^*<\frac{u^*}{1-\eta}<(1-\alpha)\omega_X<1-\alpha,
 \qquad c^*=1-u^*-i^*-m^*>0.
\]
Thus positive consumption follows from the stated conditions; it is not
an additional assumption.

For given $A_0,N_0>0$, choose the reference initial stocks as
\begin{align}
 B_0^*&=\left[
 A_0N_0(k^*)^{\alpha/[(1-\alpha)\omega_L]}(u^*)^{\omega_X/\omega_L}
 m^*\left(\frac{\chi}{g_B^*}\right)^{1/\eta}
 \right]^{\eta\omega_L/(1-\eta-\omega_X)},
 \label{eq:rewrite-uncapped-unit-initial-b}\\
 Y_0^*&=A_0N_0(k^*)^{\alpha/[(1-\alpha)\omega_L]}
 (u^*B_0^*)^{\omega_X/\omega_L},
 \qquad K_0^*=k^*Y_0^*.
 \label{eq:rewrite-uncapped-unit-initial-y-k}
\end{align}
The exponent is positive. These choices solve simultaneously the initial
production equation and
$g_B^*B_0^*=\chi(B_0^*m^*Y_0^*)^\eta$.
Set $U_0^*=u^*Y_0^*$, $M_0^*=m^*Y_0^*$, $C_0^*=c^*Y_0^*$,
$X_0^*=B_0^*U_0^*$, and
\begin{equation}
 q_0^*=\frac{M_0^*}{\eta g_B^*B_0^*}.
 \label{eq:rewrite-uncapped-unit-initial-q}
\end{equation}
Let $Y_t^*=Y_0^*e^{g_Y^*t}$, $B_t^*=B_0^*e^{g_B^*t}$,
$K_t^*=k^*Y_t^*$, $C_t^*=c^*Y_t^*$, $U_t^*=u^*Y_t^*$,
$M_t^*=m^*Y_t^*$, and $q_t^*=q_0^*e^{(g_Y^*-g_B^*)t}$.
Use the exogenous $A_t,N_t$ and set $L_t=N_t$ and $X_t^*=B_t^*U_t^*$.
The static equilibrium equations determine all remaining quantities and prices.

The two growth identities
\eqref{eq:rewrite-uncapped-unit-production-growth}
and~\eqref{eq:rewrite-uncapped-unit-research-growth} ensure that the production
equation and the RSI technology hold at every date once they hold initially.
The static monopoly condition gives $p_XX=(1-\alpha)\omega_XY$ and $U=u^*Y$.
Equation~\eqref{eq:rewrite-uncapped-unit-initial-q} gives
$M=\eta q\dot B$, which is the research first-order condition
\eqref{eq:rewrite-eq-research} with $\psi=1$.
The definitions of $k^*$, $r^*$, and $c^*$ verify the capital-return,
Euler, and resource equations. The costate equation reduces to
\[
 g_Y^*-g_B^*
 =r^*-\frac{u^*\eta g_B^*}{m^*}-\eta g_B^*.
\]
This holds by equation~\eqref{eq:rewrite-uncapped-unit-research-share}
and $(1-\eta)g_B^*=\eta g_Y^*$.
All allocations are positive and finite at every finite date.

Both transversality conditions are explicit:
\begin{align}
 e^{-\rho t}\frac{N_tK_t^*}{C_t^*}
 &=\frac{N_0k^*}{c^*}e^{-(\rho-n)t}\longrightarrow0,
 \label{eq:rewrite-uncapped-unit-household-tvc}\\
 e^{-r^*t}q_t^*B_t^*
 &=q_0^*B_0^*e^{-(\rho-n)t}\longrightarrow0.
 \label{eq:rewrite-uncapped-unit-developer-tvc}
\end{align}
Discounted developer profit is a constant multiple of
$e^{-(\rho-n)t}$. Discounted household utility is
$N_0e^{-(\rho-n)t}[\log(C_0^*/N_0)+(g_Y^*-n)t]$.
Both improper objective integrals are finite.

It remains to establish global optimality, rather than only the dated
necessary conditions. At given paths of $K,A,L$, static optimization gives
\begin{equation}
 \begin{split}
 \pi_t(B)={}&(1-\alpha)\omega_X[1-(1-\alpha)\omega_X]
 \bigl[K_t^\alpha(A_tL_t)^{(1-\alpha)\omega_L}\bigr]^{1/[1-(1-\alpha)\omega_X]}\\
 &\times\bigl[(1-\alpha)^2\omega_X^2B\bigr]^{(1-\alpha)\omega_X/[1-(1-\alpha)\omega_X]}.
 \end{split}
 \label{eq:rewrite-uncapped-unit-profit}
\end{equation}
For this verification, define the invertible state transformation and its
current-value costate by
\begin{equation}
 \nu=\frac{1-\eta}{\eta},\qquad
 S=B^\nu,\qquad p=\frac{q}{\nu B^{\nu-1}}.
 \label{eq:rewrite-uncapped-unit-transformed-state}
\end{equation}
The RSI technology becomes
\[
 \dot S=\nu\chi S^{1-\eta}M^\eta,
\]
which is jointly concave in $S,M$ for every $0<\eta<1$.
At given $K,A,L$, equation~\eqref{eq:rewrite-uncapped-unit-profit} is a
positive time-dependent coefficient times
\[
 S^{\frac{(1-\alpha)\omega_X}
 {\nu[1-(1-\alpha)\omega_X]}}.
\]
Since $\eta+\omega_X<1$ implies $(1-\alpha)\omega_X<1-\eta$, this
exponent lies strictly between zero and one. Operating profit is therefore
concave in $S$. With $p>0$, the transformed Hamiltonian
$\pi_t(S^{1/\nu})-M+p_t\nu\chi S^{1-\eta}M^\eta$ is jointly concave.
The research condition~\eqref{eq:rewrite-research-foc} and costate
equation~\eqref{eq:rewrite-costate}, with $\psi=1$, give its control
first-order condition and costate equation by the chain rule.

Write $D_t=\exp(-\int_0^t r_s\dd s)$ and
$J_T=\int_0^T D_t[\pi_t(B_t)-M_t]\dd t$.
For any admissible alternative with the same initial $B$, concavity,
the research condition, and the costate equation give, after integration
by parts,
\begin{equation}
 \widetilde J_T-J_T
 \leq-D_Tp_T(\widetilde S_T-S_T)
 \leq D_Tp_TS_T
 =\frac{D_Tq_TB_T}{\nu}\longrightarrow0.
 \label{eq:rewrite-uncapped-unit-developer-verification}
\end{equation}
The last limit follows from the developer TVC. Since the transformation
preserves feasible research paths and their payoffs, no admissible
alternative has a greater limiting payoff, or a greater upper limit of
truncated payoff, than this finite-valued candidate. The same verification
applies to nearby candidates satisfying the uncapped equations and the
developer TVC; it does not require exact balanced growth.

For the household, set $\Lambda_t=e^{-\rho t}N_t/C_t$.
Euler implies $\dot\Lambda=-r\Lambda$.
Concavity of log utility and the linear budget give, for any feasible
alternative with the same $K_0$ and the same given prices and distributions,
\[
 \int_0^T e^{-\rho t}N_t
 \left[\log(\widetilde C_t/N_t)-\log(C_t/N_t)\right]\dd t
 \leq-\Lambda_T(\widetilde K_T-K_T)
 \leq\Lambda_TK_T\longrightarrow0.
\]
This establishes global household optimality. Competitive final-good firms
face a concave constant-returns technology, and
Lemma~\ref{lem:rewrite-monopoly-choice} establishes global static monopoly
optimality. The constructed path satisfies Definition~\ref{def:rewrite-equilibrium}
with the uncapped RSI technology~\eqref{eq:rewrite-uncapped-unit-law}.

\medskip
\textit{Local convergence from nearby initial stocks.}
\phantomsection
\label{proof:rewrite-uncapped-unit-local}
Use the ordering $\xi=(\xi_K,\xi_B,\xi_C,\xi_q)'$ in
\eqref{eq:rewrite-uncapped-unit-deviations}, and let $e_K,e_B,e_C,e_q$
be the corresponding unit row vectors. Define the row vectors
\begin{align}
 y_\xi&=\left(\frac{\alpha}{1-(1-\alpha)\omega_X},
 \frac{(1-\alpha)\omega_X}{1-(1-\alpha)\omega_X},0,0\right),
 \label{eq:rewrite-uncapped-unit-output-row}\\
 m_\xi&=\frac{\eta e_B+e_q}{1-\eta},
 \label{eq:rewrite-uncapped-unit-research-row}\\
 b_\xi&=g_B^*[(\eta-1)e_B+\eta m_\xi].
 \label{eq:rewrite-uncapped-unit-efficiency-row}
\end{align}
The reduced production equation gives the log deviation of output as
$y_\xi\xi$. The research FOC implies
$M=(q\chi\eta B^\eta)^{1/(1-\eta)}$, so the log deviation of $M$ is
$m_\xi\xi$. Differentiating $g_B=\chi B^{\eta-1}M^\eta$ at the BGP
gives the level deviation $b_\xi\xi$.

The four rows of $J^*$, derived from capital accumulation, AI research,
the household Euler equation, and the developer costate equation, are
\begin{align}
 J_K^*&=\frac{1-u^*}{k^*}y_\xi-(g_Y^*+\delta)e_K
 -\frac{c^*}{k^*}e_C-\frac{m^*}{k^*}m_\xi,
 \label{eq:rewrite-uncapped-unit-jacobian-k}\\
 J_B^*&=b_\xi,
 \label{eq:rewrite-uncapped-unit-jacobian-b}\\
 J_C^*&=(r^*+\delta)(y_\xi-e_K),
 \label{eq:rewrite-uncapped-unit-jacobian-c}\\
 J_q^*&=(r^*+\delta)(y_\xi-e_K)
 -(\rho-n+\eta g_Y^*)(y_\xi-e_q-e_B)-\eta b_\xi.
 \label{eq:rewrite-uncapped-unit-jacobian-q}
\end{align}
The last row uses $U/(qB)=\rho-n+\eta g_Y^*$ at the BGP.
All entries depend on growth rates and ratios, which are independent of
$\chi$. No terminal restrictions were used to derive this Jacobian.

I next establish the two local properties
\eqref{eq:rewrite-uncapped-unit-local-spectrum}
and~\eqref{eq:rewrite-uncapped-unit-local-projection} from the parameter
restrictions, rather than impose them. Retain $\nu$ from
\eqref{eq:rewrite-uncapped-unit-transformed-state} and, for this linearization,
define
\begin{align*}
 a_K&=\frac{\alpha}{1-(1-\alpha)\omega_X},&
 a_T&=\frac{(1-\alpha)\omega_X\eta}
 {[1-(1-\alpha)\omega_X](1-\eta)},\\
 d&=\rho-n,&P&=r^*+\delta,&h&=\rho-n+\eta g_Y^*.
\end{align*}
All six quantities are positive, and the restriction $\eta+\omega_X<1$ gives
\begin{equation}
 1-a_K-a_T
 =\frac{(1-\alpha)(1-\eta-\omega_X)}
 {[1-(1-\alpha)\omega_X](1-\eta)}>0.
 \label{eq:rewrite-uncapped-unit-local-gap}
\end{equation}
Make the invertible linear change of deviations
\[
 \widehat\xi=(\xi_K,\nu\xi_B,\xi_C,\xi_q+(1-\nu)\xi_B)'.
\]
This changes neither the model nor the eigenvalues. Applying it to
\eqref{eq:rewrite-uncapped-unit-jacobian-k}--\eqref{eq:rewrite-uncapped-unit-jacobian-q}
gives a similar matrix
\begin{equation}
 \widehat J=
 \begin{pmatrix}
  \mathsf A&\mathsf B&-\mathsf C&-\mathsf V\\
  0&0&0&\mathsf D\\
  -\mathsf G&\mathsf H&0&0\\
  -\mathsf E&\mathsf F&0&d
 \end{pmatrix},
 \label{eq:rewrite-uncapped-unit-local-similarity}
\end{equation}
where the following letters abbreviate matrix coefficients only:
\begin{align*}
 \mathsf A&=d+(1-\alpha)\omega_X P,&
 \mathsf B&=\frac{(1-u^*)a_T-m^*}{k^*},&
 \mathsf C&=\frac{c^*}{k^*},\\
 \mathsf V&=\frac{m^*}{k^*(1-\eta)},&
 \mathsf D&=\frac{\eta g_Y^*}{1-\eta},&
 \mathsf G&=P(1-a_K),\\
 \mathsf H&=Pa_T,&
 \mathsf E&=\mathsf G+ha_K,&
 \mathsf F&=\mathsf H+h(1-a_T).
\end{align*}
All coefficients other than $\mathsf B$ are positive by their definitions.
The BGP allocation
ratios imply the identities
\begin{align}
 \mathsf L&\equiv\mathsf B-\frac{\mathsf Vd}{\mathsf D}
 =\frac{(1-\alpha)\omega_X\eta}{(1-\eta)k^*}>0,
 \label{eq:rewrite-uncapped-unit-local-l}\\
 \mathsf G\mathsf F-\mathsf H\mathsf E
 &=hP(1-a_K-a_T)>0.
 \label{eq:rewrite-uncapped-unit-local-minor}
\end{align}
Set $\mathsf N=\mathsf C\mathsf G+\mathsf V\mathsf E>0$ and
$\mathsf Q=\mathsf C\mathsf H+\mathsf V\mathsf F>0$.
Expanding the determinant gives the characteristic polynomial
\begin{equation}
 \begin{split}
 p(\ell)&=\det(\ell I-\widehat J)\\
 &=(\ell^2-\mathsf A\ell-\mathsf N)
 (\ell^2-d\ell-\mathsf D\mathsf F)
 -\mathsf D\mathsf E(\mathsf Q-\mathsf L\ell).
 \end{split}
 \label{eq:rewrite-uncapped-unit-local-polynomial}
\end{equation}
In particular,
\[
 p(0)=\mathsf D\mathsf C hP(1-a_K-a_T)>0.
\]
Let $\mu=(d-\sqrt{d^2+4\mathsf D\mathsf F})/2<0$, the negative
root of the second quadratic factor. Then
$p(\mu)=-\mathsf D\mathsf E(\mathsf Q-\mathsf L\mu)<0$.
Since $p(\ell)\to+\infty$ as $\ell\to-\infty$, there is one
negative root below $\mu$ and another between $\mu$ and zero.
These are the only negative roots: the cubic coefficient of $p$ is
$-(\mathsf A+d)<0$, its linear coefficient is
$\mathsf A\mathsf D\mathsf F+d\mathsf N+\mathsf L\mathsf D\mathsf E>0$,
and its constant coefficient is positive. Thus the coefficient signs of
$p(-x)$, in descending powers, are $(+,+,\text{arbitrary},-,+)$.
Descartes' rule of signs gives exactly two sign changes, regardless of
the quadratic coefficient. The two negative roots are therefore real,
distinct, and simple. Write them as $\ell_1<\mu<\ell_2<0$.
The remaining two roots have positive product and positive sum, because
$p(0)>0$ and the sum of all four roots is $\mathsf A+d>0$.
If real, both are positive; if complex, their real parts are positive.
This proves \eqref{eq:rewrite-uncapped-unit-local-spectrum}.

To establish the state projection, write a stable eigenvector in transformed
coordinates as $(x,y,z,u)'$. Its second component cannot be zero:
the second, fourth, and first rows would successively imply $u=x=z=0$.
Normalize $y=1$. The second and fourth rows then give
\[
 x(\ell)=\frac{\mathsf D\mathsf F+d\ell-\ell^2}{\mathsf D\mathsf E}.
\]
The ordering of the roots implies $x(\ell_1)<0<x(\ell_2)$.
Hence the two projected stable eigenvectors $(x(\ell_1),1)'$ and
$(x(\ell_2),1)'$ are linearly independent. Returning to $(\xi_K,\xi_B)$
only rescales their second component by $1/\nu>0$, so
\eqref{eq:rewrite-uncapped-unit-local-projection} follows as well.

The exact normalized vector field is continuously differentiable near
$\xi=0$. The spectral property just proved and the stable-manifold
theorem supply a two-dimensional local manifold
whose trajectories remain near zero and converge exponentially.
The derivative of its projection onto $(\xi_K,\xi_B)$ is $V_{KB}$.
The nonsingularity just proved and the
inverse-function theorem therefore make this manifold locally the
graph of a continuously differentiable function selecting $(\xi_C,\xi_q)$
from the two initial stock deviations. Since exponentiation is a local
diffeomorphism, the same holds for $(C_0,q_0)$ as functions of $(K_0,B_0)$.
More explicitly, for some constants $a,H>0$, trajectories on a sufficiently
small stable manifold satisfy
\[
 \|\xi(t)\|\leq H e^{-at}\|\xi(0)\|,\qquad t\geq0.
\]
For any sufficiently small open neighborhood $\mathcal V$ of zero,
choose a ball contained in $\mathcal V$. Continuity of the graph selecting
the two jump variables permits an open neighborhood $\mathcal U$ of the
initial stock deviations for which $H\|\xi(0)\|$ is smaller than that
ball's radius. Every initial stock pair in $\mathcal U$ therefore has a
trajectory staying in $\mathcal V$ for all $t\geq0$ and converging to zero.
Exponentiation maps $\mathcal U$ to an open neighborhood of
$(K_0^*,B_0^*)$. Shrinking these neighborhoods preserves positivity,
feasibility, and existence for every $t\geq0$. Local uniqueness follows
because every trajectory remaining in this local region and converging
to zero belongs to the same stable manifold, which is a graph over stocks.
This does not assert convergence from arbitrary initial $C_0,q_0$ or
exclude equilibrium trajectories that leave the region.

The statement also permits nearby $A_0,N_0$. Equations
\eqref{eq:rewrite-uncapped-unit-initial-b}
and~\eqref{eq:rewrite-uncapped-unit-initial-y-k} make the reference
$B_0^*,K_0^*$ smooth functions of these two positive initial values.
After normalization by the corresponding BGP, the exact vector field
depends only on the parameters and BGP ratios, not on $A_0,N_0$.
The set of $(A_0,N_0,K_0,B_0)$ whose stock deviations belong to
$\mathcal U$ is therefore open and contains the reference initial condition.
From each such initial condition, the construction converges to the BGP
associated with its own $A_0,N_0$, not necessarily to the same levels as
the original reference path.

Since $f(\xi(t))\to0$, growth rates converge to their BGP values, not only
normalized levels. The static equations give the same convergence for
the interest rate, wage growth, and allocation ratios. The unit-elastic
labor income share remains $(1-\alpha)\omega_L$ at every date.

Exponential convergence implies that $r_t-r^*$ is integrable. Thus
$D_t/e^{-r^*t}$ tends to a finite positive constant, and both TVC expressions
equal their BGP counterparts in
\eqref{eq:rewrite-uncapped-unit-household-tvc}
and~\eqref{eq:rewrite-uncapped-unit-developer-tvc} times factors tending
to finite positive limits. Both TVCs hold. Output and all expenditure
flows are bounded multiples of $e^{g_Y^*t}$, while log consumption per
person is a linear function of time plus a bounded deviation. Both
objective integrals therefore remain finite. The global concavity
verification established above
applies at these new price and allocation paths, proving agent optimality.
Static optimality and market clearing then establish equilibrium, not
merely a solution of the necessary conditions.
\end{proof}
